%% CADENCE — unconstrained recurrence intensities for TKG forecasting
%% IEEEtran (ICDE / general IEEE conference). For IJCAI switch to ijcai.sty.
\documentclass[conference]{IEEEtran}
\IEEEoverridecommandlockouts

\usepackage{amsmath,amssymb,amsthm}
\usepackage{graphicx}
\usepackage{booktabs}
\usepackage{multirow}
\usepackage{url}
\usepackage[hidelinks]{hyperref}
\usepackage{xcolor}
%% Overleaf has Inkscape installed, so \includesvg works there directly and the
%% figure stays vector. If you compile locally without Inkscape, either run
%%   inkscape architecture.svg --export-filename=architecture.pdf
%% and swap \includesvg for \includegraphics, or compile with -shell-escape.
\usepackage{svg}
\svgsetup{inkscapelatex=false}

\newtheorem{proposition}{Proposition}
\newcommand{\todo}[1]{\textcolor{red}{[#1]}}
\newcommand{\pending}[1]{\textcolor{blue}{#1$^{\dagger}$}}

\begin{document}

\title{Beyond Monotone Decay: Unconstrained Recurrence\\
Intensities for Temporal Knowledge Graph Forecasting}

\author{\IEEEauthorblockN{Anonymous}
\IEEEauthorblockA{\textit{Institution} \\ \textit{Address} \\ email}}

\maketitle

\begin{abstract}
Recurrence dominates temporal knowledge graph (TKG) forecasting: a plain
recurrency baseline matches RE-GCN, CyGNet, TiRGN and L2-TKG on the standard
benchmarks. Every recurrence mechanism in this literature scores a historical
candidate with $f(\mathrm{count})\cdot g(\Delta t)$ for some non-increasing
$g$, including the exponential kernel of the recent neural-Hawkes model
GAttNHP, which learns the decay rate but keeps the form. We show that this
family is provably unable to rank a large and directly measurable share of
queries: whenever a distractor has both smaller elapsed time and larger count
than the answer, no non-increasing $g$ and no non-decreasing $f$ can place the
answer first. We call such queries \emph{monotone-blocked}, note that the
property is defined without reference to any model, and measure it at
$41.6\%$ of YAGO, $29.3\%$ of WIKI, $26.6\%$ of ICEWS18 and $15.8\%$ of GDELT
test queries. We then introduce \textsc{Cadence}, which abandons the
product form altogether: its recurrence branch is an unconstrained function of
thirteen inter-arrival statistics of a candidate's observed history, so it is
neither monotone in $\Delta t$ nor separable in the count, and it recovers the
classical exponential kernel as one point in its hypothesis space. The two
branches are combined as what they are --- intensities of superposed point
processes --- through a log-sum-exp rather than a learned gate, which removes
the branch collapse we observed with gating. Restricting the branch to the
published monotone form while holding everything else fixed costs $1.83$,
$8.11$, $5.39$ and $4.57$ MRR on the four benchmarks, so the functional form,
not the backbone, carries the gain. \textsc{Cadence} exceeds the published
state of the art on all four benchmarks under the time-aware filtered
protocol, which we verify against the reference implementation.
\end{abstract}

\begin{IEEEkeywords}
temporal knowledge graphs, link forecasting, temporal point processes,
recurrence
\end{IEEEkeywords}

\section{Introduction}
\label{sec:intro}

\todo{Draft. The argument, in order:}
\begin{enumerate}
  \item Recurrence dominates TKG forecasting, and the field knows it: a plain
        recurrency baseline matches four published architectures.
  \item So the question is not whether to model recurrence, but what existing
        recurrence models cannot express.
  \item They share one functional form, and that form has a blind spot that
        is provable and whose size is measurable from the data alone.
  \item The blind spot covers $15.8$--$41.6\%$ of test queries, so it is not a
        corner case that a structural branch can be assumed to absorb.
  \item Contributions: Proposition~\ref{prop:blocked} and the model-free
        diagnostic it induces; the \textsc{Cadence} recurrence intensity, which
        strictly contains the family the proposition constrains; a
        superposition rule that follows from the point-process reading and
        admits no gate; results on four benchmarks under a protocol verified
        against the reference implementation; and a stratified evaluation that
        tests the mechanism rather than only the headline metric.
\end{enumerate}

\section{Related Work}
\label{sec:related}

\paragraph{Recurrence and copy mechanisms}
CyGNet, CENET, TiRGN, RE-GCN and DaeMon all score historical candidates by
frequency weighted by exponential decay; CENET adds an explicit
historical/non-historical oracle trained with a contrastive objective. All are
instances of $f(\mathrm{count})\cdot g(\Delta t)$.

\paragraph{Point processes for TKG}
GAttNHP casts TKG extrapolation as a neural Hawkes process with base, self-
and group-excitation, but its kernel is $\exp(-\gamma_u \Delta t)$ with a
learned per-entity rate. It is therefore inside the family constrained by
Proposition~\ref{prop:blocked}. Our contribution is orthogonal to the
point-process framing itself and concerns the shape of the kernel.

\paragraph{Structural encoders}
RE-GCN evolves entity representations over snapshots with an R-GCN and a GRU.
DiMNet adds cross-time perception at equal hop depth and disentangles node
features into active and stable factors; its ablation attributes $-11.38$ MRR
to disentanglement, $-9.62$ to virtual-subgraph refinement and $-4.97$ to
multi-span. \emph{We adopt multi-span evolution and active/stable
disentanglement as backbone and claim no novelty for them.}

\paragraph{Path-based reasoning}
DaeMon carries an NBFNet-style query-conditioned state across snapshots with a
time-aware gate and contains no entity embeddings, so a never-seen candidate
is still scored on path evidence. We reimplement this as an optional third
intensity and report it in Section~\ref{sec:negative} as a negative result at
our scale.

\paragraph{Evaluation protocol}
Published numbers use three incompatible settings. RE-GCN, TiRGN, CENET,
DaeMon, TiPNN and DiMNet report time-aware filtered; GAttNHP reports raw. We
report all three and compare only within one.

\section{Preliminaries}
\label{sec:prelim}

A TKG is a set of quadruples $(s,r,o,t)$. Forecasting asks for $o$ given
$(s,r,?,t)$ using only facts strictly before $t$. We report MRR and
Hits@$\{1,3,10\}$. Ties are resolved to their average position,
\begin{equation}
\mathrm{rank} = 1 + \#\{\text{strictly better}\}
              + \tfrac{1}{2}\bigl(\#\{\text{tied}\} - 1\bigr),
\end{equation}
because counting only strictly-better scores inflates every metric for models
that assign many candidates an identical score. Under that convention an early
copy-only variant of our own model reported $\mathrm{H@}10 = 99.93$ on YAGO,
against $93.34$ for DaeMon.

\section{The Blind Spot of Monotone Recurrence}
\label{sec:blocked}

\subsection{Proposition}

\begin{proposition}
\label{prop:blocked}
Let a scorer assign $s(o) = f(c_o)\,g(\Delta t_o)$ with $f$ non-decreasing,
$g$ positive and non-increasing. Let $o^\star$ be the answer and $o$ a
distractor with $\Delta t_o \le \Delta t_{o^\star}$ and
$c_o \ge c_{o^\star}$. Then $s(o) \ge s(o^\star)$.
\end{proposition}

\begin{proof}
$f(c_o) \ge f(c_{o^\star})$ since $f$ is non-decreasing, and
$g(\Delta t_o) \ge g(\Delta t_{o^\star})$ since $g$ is non-increasing. Both
factors are positive, so the product is at least as large.
\end{proof}

No choice of decay rate --- fixed, per-relation, or learned per entity as in
GAttNHP --- and no reweighting of counts ranks $o^\star$ strictly first. We
call such a query \emph{monotone-blocked}; the property involves only the
query and its candidate history, and no model.

\subsection{The rate is large}

\begin{table}[t]
\centering
\caption{Monotone-blocked rate, measured on 20{,}000 sampled test queries per
dataset. No model is involved.}
\label{tab:blocked}
\begin{tabular}{lccc}
\toprule
dataset & answer has history & blocked (of those) & blocked / all \\
\midrule
YAGO    & 92.8\% & 44.8\% & \textbf{41.6\%} \\
WIKI    & 87.0\% & 33.7\% & \textbf{29.3\%} \\
ICEWS18 & 47.7\% & 55.7\% & \textbf{26.6\%} \\
GDELT   & 24.9\% & 63.4\% & \textbf{15.8\%} \\
\bottomrule
\end{tabular}
\end{table}

\paragraph{Scope of the bound}
Proposition~\ref{prop:blocked} constrains a recurrence scorer \emph{in
isolation}. Published systems pair such a term with a structural branch that
is not subject to it and can rescue a blocked query. What
Table~\ref{tab:blocked} establishes is that on $15.8$--$41.6\%$ of queries the
recurrence component cannot produce the correct ordering and must delegate.

\subsection{Two regimes}

Table~\ref{tab:blocked} separates the benchmarks into two regimes, and the rest
of the paper is read against that split. On YAGO and WIKI almost every answer
has a history ($92.8\%$, $87.0\%$), so recurrence is where the task is decided
and the blocked share of \emph{all} queries is correspondingly large. On
ICEWS18 and GDELT most answers have no history at all ($47.7\%$ and $24.9\%$
do), so the blocked share of all queries is smaller --- but conditional on
having a history it is \emph{higher}, $55.7\%$ and $63.4\%$. The blind spot is
therefore not confined to the slowly-evolving datasets; on the event datasets
it is denser, over a smaller pool. Section~\ref{sec:strat} evaluates on these
strata separately rather than reporting only the pooled metric.

\section{\textsc{Cadence}}
\label{sec:method}

\begin{figure*}[t]
\centering
\includesvg[width=\textwidth]{architecture}
\caption{Overview of the model. \emph{Left}: the structural branch. Each
historical snapshot $\mathcal{G}(t_i)$ is passed through three multi-span GCN
layers whose message at depth $\ell$ also reads depth $\ell$ at $t_{i-1}$, is
disentangled into an active and a stable factor, and is carried across time by a
GRU, giving the evolved entity table $E(t)$. A ConvTransE decoder scores every
entity against $[E(t)[s]\,\|\,R[r]]$. \emph{Right}: the recurrence branch, drawn
in full because it is where the contribution lies. For each candidate $o$ we read
its observed occurrence history off a precomputed CSR index and summarise it by
thirteen inter-arrival statistics, which an unconstrained MLP maps to a scalar.
The two timelines $o^{\star}$ and $o'$ are the configuration of
Proposition~\ref{prop:blocked}: the distractor has both the smaller $\Delta t$
and the larger count, so every score of the form $f(\text{count})\cdot g(\Delta
t)$ with $g$ non-increasing ranks it at least as high as the answer, whatever
decay rate is fitted. \emph{Bottom}: the branches are point-process intensities,
so they are superposed by $\log\lambda = \mathrm{logaddexp}(f_{\mathrm{struct}},
f_{\mathrm{rec}})$ on the candidate set and by $f_{\mathrm{struct}}$ elsewhere.
No mixing weight is learned, so there is no gate that can collapse onto one
branch.}
\label{fig:arch}
\end{figure*}

\subsection{Superposition, not gating}

Fig.~\ref{fig:arch} gives the whole model. A query is a draw from a
superposition of two marked point processes on the entity space,
\begin{equation}
\lambda(o) = \lambda_{\mathrm{struct}}(o \mid \mathcal{G}_{<t}, s, r)
           + \lambda_{\mathrm{rec}}(o \mid \mathcal{H}_o, s, r).
\end{equation}
For a superposition the probability that the next mark is $o$ is
$\lambda(o)/\sum_{o'}\lambda(o')$, so a softmax cross-entropy on
\begin{equation}
\mathrm{score}(o) = \log \lambda(o) =
\begin{cases}
\mathrm{logaddexp}\bigl(f_{\mathrm{struct}}(o), f_{\mathrm{rec}}(o)\bigr) & o \in S,\\
f_{\mathrm{struct}}(o) & \text{otherwise},
\end{cases}
\end{equation}
is the correct discrete-choice likelihood. Two consequences matter in
practice. Superposition \emph{adds intensities}, so the branches combine
through $\mathrm{logaddexp}$ and not through a sum of logits; a sum of logits
multiplies intensities, which has no point-process meaning and allows one
branch to rescale the other. And there is no mixing gate, hence no gate to
collapse: earlier variants of this work with a learned scalar or per-query
gate collapsed to a single branch within a few epochs.

\subsection{The recurrence intensity}

The proposition constrains a functional \emph{form}, not a parameterisation, so
no amount of fitting inside that form escapes it. We therefore leave the form.
For a candidate $o$ we summarise its observed history by a vector
$\phi_o \in \mathbb{R}^{13}$ and set
\begin{equation}
\log \lambda_{\mathrm{rec}}(o) =
\mathrm{MLP}\bigl[\,\mathrm{LN}(\phi_o) \,\|\, u(r) \,\|\, v(o)\,\bigr] + b,
\end{equation}
where $u$ and $v$ are learned relation and entity contexts, the trunk is two
Linear--LayerNorm--GELU blocks, and $b$ is a scalar initialised at $-2$.

The thirteen statistics are the count, the elapsed time $\Delta t$ and the mean
inter-arrival gap of each of the $(s,r,o)$, $(s,\cdot,o)$ and $(r,\cdot,o)$
histories, together with the dispersion of the gaps, the span and rate of the
$(s,r,o)$ history, and a binary has-history indicator. They are read from a
compressed sparse-row temporal index built once per dataset and carry no
gradient, so the branch adds no message passing to the backbone.

Two properties matter. First, the family in
Proposition~\ref{prop:blocked} is contained in this one: taking the MLP to
compute $\log a + p\log(1+c) - b\,\Delta t$ recovers the classical
count-times-decay scorer exactly, so we cannot lose by generalising, and
Section~\ref{sec:mech} measures what we gain by doing so. Second, nothing
forces the score to decrease in $\Delta t$ or to increase in the count, and it
is a joint function of both together with seven further statistics --- so the
blocked configuration of Proposition~\ref{prop:blocked} is reachable, which is
precisely what the product form forbids.

\subsection{Structural intensity}

The backbone is adopted from prior work: a composition GCN over each snapshot
with a cross-time link at equal layer depth, node states split into
complementary active and stable factors with a drift penalty on the stable
factor across adjacent snapshots, a GRU across snapshots, and a ConvTransE
decoder over the full entity table. We use LayerNorm rather than BatchNorm in
the decoder: batches here are snapshots, so batch size is however many facts
occurred at one timestamp and varies from $72$ to $4{,}650$ across our
datasets, and once a long snapshot is split into query chunks the result would
otherwise depend on the split.

\section{Experimental Setup}
\label{sec:setup}

Four benchmarks, time-aware filtered protocol, baselines quoted from the
DaeMon table (YAGO, WIKI) and the DiMNet table (ICEWS18, GDELT), both of which
report that protocol. $d=200$, three GCN layers, history of ten snapshots
(five on GDELT), $5.96$M parameters. One A100 per run: YAGO 42\,min, ICEWS18
1\,h\,15, WIKI 3\,h, GDELT 8\,h.

\subsection{The protocol is verified against the reference}
\label{sec:protocol}

A subtle difference in filtering would invalidate margins of the size we
report, so we checked our harness against RE-GCN's \texttt{rgcn/utils.py},
which is the de-facto standard for this benchmark family. The reference
appends inverse triples in \texttt{predict()} and scores both directions; so
do we, and the scored query counts are identical on all four datasets. The
reference sets every other true answer at the query timestamp to $-10^7$ and
restores the target; we do the same with $-\infty$. The reference builds the
filter from the test split alone whereas we build it from all splits; the
filter is keyed by timestamp, so the two coincide when splits are temporally
disjoint. We verified disjointness on the actual files --- zero overlapping
timestamps on all four datasets --- and then verified directly that the two
answer sets agree on $4{,}000$ sampled test queries per dataset, with zero
differences.

One difference remains. The reference ranks by sort position, so a target tied
with $k-1$ others lands anywhere in that group; we take the average position,
which is the expectation of that draw. Ours is therefore unbiased with respect
to the reference and cannot be optimistic.

\section{Results}
\label{sec:results}

\begin{table}[t]
\centering
\caption{Main results, time-aware filtered, $\times 100$, mean $\pm$ std over
seeds. $^{\dagger}$ single seed, pending.}
\label{tab:main}
\begin{tabular}{llcccc}
\toprule
dataset & method & MRR & H@1 & H@3 & H@10 \\
\midrule
\multirow{2}{*}{YAGO}
 & DaeMon  & 91.59 & 90.03 & \textbf{93.00} & \textbf{93.34} \\
 & \textsc{Cadence} & \textbf{91.64}\,$\pm$0.04 & \textbf{90.32}\,$\pm$0.08 & 92.88 & 93.08 \\
\midrule
\multirow{2}{*}{WIKI}
 & DaeMon  & 82.38 & 78.26 & \textbf{86.03} & \textbf{88.01} \\
 & \textsc{Cadence} & \textbf{82.94}\,$\pm$0.15 & \textbf{79.84}\,$\pm$0.28 & 85.77 & 86.95 \\
\midrule
\multirow{2}{*}{ICEWS18}
 & DiMNet  & 34.13 & 23.29 & 38.42 & \textbf{55.80} \\
 & \textsc{Cadence} & \textbf{34.59}\,$\pm$0.14 & \textbf{24.78}\,$\pm$0.11 & \textbf{39.06} & 53.55 \\
\midrule
\multirow{2}{*}{GDELT}
 & DiMNet  & 21.93 & 14.03 & 23.57 & 37.49 \\
 & \textsc{Cadence}\pending{} & \textbf{26.34} & \textbf{17.59} & \textbf{29.07} & \textbf{43.44} \\
\bottomrule
\end{tabular}
\end{table}

Margins over the state of the art relative to the seed standard deviation are
$13.5\times$ on ICEWS18 H@1, $5.6\times$ on WIKI and $3.6\times$ on YAGO.
We lose H@10 on three of four datasets; Section~\ref{sec:strat} shows why.

\subsection{The mechanism ablation}
\label{sec:mech}

The claim is about a functional form, so the ablation replaces the form and
nothing else. We restrict the recurrence branch to exactly the published
scorer, $\log a(r,o) + p(r,o)\log(1+c) - b(r,o)\,\Delta t$ with $p, b \ge 0$
predicted from $(r,o)$ alone, and hold the backbone, the superposition rule,
the optimiser and the schedule fixed. A unit test sweeps $\Delta t$ and $c$ over
twenty seeds and asserts that the restricted branch is monotone in both and
that the full branch is not.

\begin{table}[t]
\centering
\caption{Full model against the monotone kernel, MRR.}
\label{tab:mono}
\begin{tabular}{lccc}
\toprule
dataset & full & monotone kernel & $\Delta$ \\
\midrule
YAGO    & 91.64 & 89.81 & $+1.83$ \\
WIKI    & 82.94 & 74.83 & $+8.11$ \\
ICEWS18 & 34.59 & 29.20 & $+5.39$ \\
GDELT   & 26.34 & 21.77 & $+4.57$ \\
\bottomrule
\end{tabular}
\end{table}

The monotone variant is itself near the state of the art on GDELT ($21.77$
against $21.93$): the backbone is competitive and the kernel adds on top.

\subsection{Where the gain lands}
\label{sec:strat}

A pooled metric cannot say whether a gain came from the mechanism this paper
argues for. We therefore partition the test set into three disjoint strata:
\emph{blocked}, where a distractor dominates the answer in both $\Delta t$ and
count; \emph{clean}, where the answer has a history and is not dominated; and
\emph{no-history}, where the answer never occurred with $(s,r)$ before. Only
the first is the target of Proposition~\ref{prop:blocked}.

\todo{Insert the stratified table once the remaining seeds land. The finding so
far, as H@1 gained over the monotone kernel: on YAGO the blocked stratum gains
$+5.66$ against $+1.77$ on clean, and on WIKI $+18.31$ against $+6.61$ --- the
gain concentrates where the proposition says the restricted form cannot
succeed. On ICEWS18 and GDELT the ordering inverts, which we attribute to those
datasets being dominated by the no-history stratum, where neither form applies;
state this, do not hide it.}

\section{Two Negative Results}
\label{sec:negative}

\paragraph{Deep supervision on the structural branch hurts}
Paired on three seeds, adding a cross-entropy on the structural scores alone
costs $0.87$ MRR and $1.66$ H@1 on YAGO, in the same direction on every seed,
and raises the seed standard deviation from $0.04$ to $0.42$. It was
introduced to stop recurrence from starving the structural branch; on a
$92.8\%$-recurrent dataset that trade is not worth making.

\paragraph{A path branch does not pay at our scale}
We reimplemented DaeMon's query-conditioned propagation as a third intensity,
with no entity embeddings. It doubled H@1 on the stratum where the answer has
no $(s,r)$ history ($0.79 \rightarrow 1.68$ on YAGO) --- exactly the stratum
it targets --- and still lost overall ($91.46$ against $91.64$ MRR), because
that stratum is only $7.3\%$ of YAGO. Its state is
(queries $\times$ entities $\times$ dim), roughly $50\times$ the arithmetic of
the other branches.

\section{Limitations}

The stratum where the answer has no $(s,r)$ history is the largest single pool
of failures on the event datasets: $50.1\%$ of ICEWS18 at $5.65$ H@1 and
$40.4\%$ of GDELT at $1.63$. The recurrence branch is silent there by
construction and the structural branch is weak; we do not close this.
H@10 is below the state of the art on three datasets. Baselines are quoted
rather than rerun, although Section~\ref{sec:protocol} removes the filtering
and query-count risk. We report three seeds, and one on GDELT.

\section*{Acknowledgements}
\todo{Per venue policy, disclose tooling used in preparation if required.}

\bibliographystyle{IEEEtran}
\bibliography{refs}

\end{document}
